\subsection{Modules over a Simple Ring}

Lemma 3. --- Let A be a simple ring, and let S be a simple A-module. Denote the opposite field of the commutant of S by D. Every A-module is isomorphic to an A-module of the form \(S \otimes_{D} V\) , where V is a left vector space over the field D.

This follows from Lemma 1 of VIII, p.~120 and Morita's theorem (VIII, p.~103).

PROPOSITION 2. --- Let A be a simple ring and S a simple A-module.

\begin{enumerate}
\def\labelenumi{\alph{enumi})}
\item
  Every A-module is projective and isotypical of type S, hence semisimple. If it has length a, then it is isomorphic to \(\mathrm{S}^{(a)}\) .
\item
  Every nonzero A-module is generating.
\item
  Two A-modules are isomorphic if and only if they have the same length.
\end{enumerate}

Denote the opposite field of the commutant of S by D. By Lemma 1 of VIII, p.~120, S is an invertible (A,D)-bimodule. Let M be an A-module; by Lemma 3, it is isomorphic to a module of the form \(S \otimes_{D} V\) , where V is a left vector space over the field D.

The D-module V is projective and isotypical of type \(D_{s}\) ; it is generating if and only if it is not reduced to 0. Finally, the length of \(S \otimes_{D} V\) is equal to the dimension of the D-vector space V, and two vector spaces are isomorphic if and only if they have the same dimension. Proposition 2 now follows in view of Propositions 10 of VIII, p.~109 and 12 of VIII, p.~110.

Let \(r \geqslant 1\) be an integer. We say that a cardinal \(\mathfrak{a}\) is divisible by \(r\) if there exists a cardinal \(\mathfrak{b}\) such that \(\mathfrak{a} = r\mathfrak{b}\) . This is the case if \(\mathfrak{a}\) is infinite because we have \(r\mathfrak{a} = \mathfrak{a}\) (Set Theory, III, §6, No.~3, p.~188, Corollary 3). It follows from this remark that if the cardinal \(\mathfrak{a}\) is divisible by \(r\) , then there exists a unique cardinal \(\mathfrak{b}\) such that \(\mathfrak{a} = r\mathfrak{b}\) .

COROLLARY. --- Let k be a commutative field, and let A be a simple k-algebra of finite degree over k. Every simple A-module is finite-dimensional over k. Two A-modules are isomorphic if and only if their dimensions over k are equal.

Every simple A-module is isomorphic to a quotient of \(A_{s}\) , hence finite-dimensional over k. The corollary then follows from Proposition 2, c).

PROPOSITION 3. --- Let A be a simple ring. An A-module M is free if and only if its length is divisible by the length of A. If that is the case, then all bases of M have the same cardinal, denoted by \(\dim_{\mathrm{A}}(\mathrm{M})\) (II, §7, No.~3, p.~294, Remark 2) and determined by the relation

\[
\mathrm{long} _ {\mathrm{A}} (\mathrm{M}) = \mathrm{long} (\mathrm{A}) \cdot \mathrm{dim} _ {\mathrm{A}} (\mathrm{M}).\tag{1}
\]

Suppose that M is free, and let \((e_{i})_{i\in\mathrm{I}}\) be a basis of M. The A-module M is the direct sum of the A-modules \(Ae_{i}\) , which are each isomorphic to \(A_{s}\) . Set \(r=\operatorname{long}_{\mathrm{A}}(\mathrm{A}_{s})\) ; this is an integer greater than or equal to 1 (VIII, p.~119, Proposition 1). We have \(\operatorname{long}_{\mathrm{A}}(\mathrm{M})=r\operatorname{Card}(\mathrm{I})\) by formula (13) of VIII, p.~73.

Conversely, suppose that the cardinal \(\operatorname{long}_{\mathrm{A}}(\mathrm{M})\) is divisible by r. Let a be the cardinal such that \(\operatorname{long}_{\mathrm{A}}(\mathrm{M}) = r\mathfrak{a}\) . Then the A-module M has the same length as \(\mathrm{A}_{s}^{(\mathfrak{a})}\) , hence is isomorphic to it by Proposition 2. This proves that M is free.

PROPOSITION 4. --- Let A be a simple ring and M a nonzero A-module. Denote the endomorphism ring of the A-module M by B, and view M as a left B-module.

\begin{enumerate}
\def\labelenumi{\alph{enumi})}
\item
  The mapping \(a \mapsto a_{M}\) is an isomorphism from A to the endomorphism ring of the B-module M.
\item
  Suppose that M has finite length as an A-module. Then the ring B is simple, and we have
\end{enumerate}

\[
\mathrm{long} _ {\mathrm{A}} (\mathrm{M}) = \mathrm{long} (\mathrm{B}) \text{ and } \mathrm{long} _ {\mathrm{B}} (\mathrm{M}) = \mathrm{long} (\mathrm{A}).\tag{2}
\]

The A-module M is generating by Proposition 2 of VIII, p.~122. By definition, we have \(\mathrm{B} = \mathrm{A}_{\mathrm{M}}'\) , and assertion a) therefore follows from Theorem 2 of VIII, p.~82.

Suppose that M is an A-module of finite length. Choose a simple A-module S, and let D be the opposite field of the endomorphism ring of S. By Lemma 3, the A-module M is isomorphic to a module of the form \(S \otimes_{D} V\) , where V is a left vector space over the field D. The vector space V is finite-dimensional. By Theorem 3 of VIII, p.~64, the ring B is isomorphic to \(\operatorname{End}_{\mathrm{D}}(\mathrm{V})\) ; by Lemma 2 (VIII, p.~121), the ring B is therefore simple. Taking Remark 1 of VIII, p.~63 into account, we obtain the equalities

\[
\mathrm{long(B)} = \mathrm {dim_ {D} (V)} = \mathrm {long_ {A} (M)}.
\]

By Remarks 1 of VIII, p.~63 and 3 of VIII, p.~121, we have the relations

\[
\mathrm{long} _ {\mathrm{B}} (\mathrm{M}) = \mathrm{long} _ {\mathrm{End} _ {\mathrm{D}} (\mathrm{V})} (\mathrm{S} \otimes_ {\mathrm{D}} \mathrm{V}) = \dim_ {\mathrm{D}} (\mathrm{S}) = \mathrm{long} (\mathrm{A}),
\]

which gives the last relation.

